% Draft based on the screened reading in related_works/PAPER_REVIEW.md.
% This section states mechanisms and evaluation limits, not measured DACO gains.
\label{sec:related-works}

Energy-efficient clustering in underwater wireless sensor networks involves
three coupled decisions: selecting cluster heads (CHs), assigning their members,
and forwarding collected data to the sink. A CH that serves a dense neighborhood
may reduce collection distances, yet become a bottleneck when it also forwards
traffic from other clusters. Related studies address these decisions through
energy-aware selection, unequal clustering, routing constraints, and data
reduction.

\subsection{Energy-aware cluster-head selection}

Xiao and Huang~\cite{Xiao2020ACOCR} proposed ACOCR, which partitions a
three-dimensional underwater deployment into cubes and selects CHs using
residual energy, sink distance, and average distance to other nodes in each
cube. An improved ant colony optimization (ACO) algorithm then searches for
inter-cluster routes. Its fixed spatial partition and static-node assumptions
limit how its conclusions transfer to uneven or mobile deployments; the
authors also leave underwater multipath effects outside their evaluation.
The distinction between CH selection and route search is important: ACOCR
uses ACO to construct routes after the CHs have been selected.

Other schemes use different selection and collection mechanisms.
EOCA~\cite{Yu2020EOCA} combines neighborhood coverage, residual energy, and
sink distance with an energy-dependent communication range. Its evaluation
includes cooperative communication and autonomous underwater vehicle (AUV)
assistance, while the collection trajectory is outside scope.
EECMR~\cite{Nguyen2021EECMR} instead uses an explicit CH-selection rule based
on energy, sink distance, and role duration, with relays connecting depth
layers. Its performance depends on communication range and incurs
re-clustering costs as topology changes.
DEKCS~\cite{Omeke2021DEKCS} combines distance- and energy-constrained
$k$-means clustering with adaptive CH thresholds and AUV collection.
However, its experiments use short-range underwater radio-frequency links;
they do not directly validate performance in an acoustic network. The authors
also acknowledge that favoring dense regions can reduce fairness to outliers.

\subsection{Density, unequal membership, and forwarding burden}

Density-based clustering and energy balancing have also been studied in
terrestrial networks. IS-$k$-means~\cite{Zhu2021SoftKMeans} initializes
clusters using density peaks and kernel density estimation, reassigns
boundary nodes, and rotates CH duties among multiple candidates. Its
single-hop radio model differs from underwater multi-hop forwarding.
PUDCRP~\cite{Ruan2019PUDCRP} uses particle swarm optimization (PSO) to
place unequal circular clustering regions according to node distribution.
Smaller clusters near the sink reserve energy for relaying, but its evaluation
assumes an ideal medium access control layer and omits control-message
energy costs. These studies show that local density is useful, while also
demonstrating why cluster geometry alone does not characterize actual
forwarding consumption.

UCPSO~\cite{Hou2022UCPSO} explicitly accounts for this additional burden in
underwater acoustic networks. It searches for CH sets using discrete PSO and
adapts cluster sizes using residual energy, sink distance, and inter-cluster
forwarding participation. Thus, a CH carrying more relay traffic can serve
fewer local members. The authors nevertheless identify connectivity and
transmission reliability as unexamined aspects of their energy-focused
evaluation. GTC~\cite{Xing2021GTC} provides another balancing mechanism:
it models CH participation as a game and uses incentives to coordinate
selection and rotation within nonuniform layers. Its static, reliable-link
assumptions and sink-side coordination define the scope of its conclusions.

\subsection{Routing feasibility, capacity, and information preservation}

Energy balance must be considered alongside feasible data delivery.
IEBR~\cite{Feng2019IEBR} uses depth thresholds and energy-level partitions
to change relay choices while constraining transmission distance and link
quality. LLCP~\cite{Mei2024LLCP} selects cluster membership according to
node load, location, and intra-cluster transmission capacity. Avoiding
congestion may require more clusters; its evaluated objective is delay and
structural complexity rather than battery lifetime.
For mobile deployments, QLAR~\cite{Li2025QLAR} incorporates predicted
link lifetime into Q-learning-based next-hop selection. Its adaptive neighbor
updates expose a further trade-off: frequent control messages improve
freshness but increase energy consumption and collisions.

Some energy savings also depend on changing the transmitted information.
The two-tier PSO scheme of Chen et al.~\cite{Chen2020TwoTierPSO} optimizes
clustering and routes for Bayesian compressive sensing, balancing consumption
with signal reconstruction quality. RTC~\cite{Ahmed2019RTC} suppresses
statistically similar readings through a cluster/region hierarchy, but
aggressive reduction can discard original information. Similarly,
EECRP~\cite{Xiao2021EECRP} combines neural-network data fusion with
genetic-algorithm route search, introducing training and optimization costs
beyond simple CH-selection rules. Such gains cannot be attributed solely
to clustering when the delivered information or collector assistance differs.

\subsection{Positioning of DACO}

These studies motivate distinguishing structural density from the energy
burden of a completed routing tree. DACO investigates a combination of
three-dimensional grid-coverage density, evaluated relative node depletion,
and live node-pheromone feedback for CH-set selection. Its density-dependent
penalty uses the consumption predicted after constructing and evaluating a
candidate tree, rather than assuming that local density alone measures relay
load. The relevant comparison is therefore whether this combination improves
the energy, depletion-balance, and computation trade-off under common
acoustic, routing, and data-delivery assumptions. Comparisons with AUV-assisted,
compressed-sensing, or duty-cycled schemes must account for their different
service conditions. No universal performance or novelty claim follows from
using ACO, density information, or energy balancing in isolation.
